\begin{helper}
For every \(\eta>0\) there is \(n_0\) such that, for every \(n>n_0\), if
\[
L=\log n,\qquad k_R=(1+\eta)\sqrt{nL},
\]
then
\[
1<L
\qquad\text{and}\qquad
16L^2\sqrt L\le k_R.
\]
\end{helper}
\begin{proof}
Since \(L=\log n\to\infty\), we may first enlarge the threshold so that
\(1<L\).  Put \(\delta=1/8\).  The standard logarithm-power comparison
\(\log n=o(n^{1/5})\) lets us further enlarge the threshold so that
\[
L\le \delta n^{1/5}.
\]
All quantities are then nonnegative, so raising both sides to the power
\(5/2\) gives
\[
L^{5/2}\le \delta^{5/2}n^{1/2}=\delta^{5/2}\sqrt n.
\]
Because \(\delta\le1\), we have \(\delta^{5/2}\le\delta^2=1/64\), and hence
\[
16L^{5/2}\le \sqrt n.
\]
The condition \(L>1\) implies \(n\le nL\), so \(\sqrt n\le\sqrt{nL}\).  Since
\(1+\eta\ge1\), this yields
\[
16L^2\sqrt L=16L^{5/2}
\le (1+\eta)\sqrt{nL}=k_R.
\]
\end{proof}
